% Appendix: standalone or \input into main paper
% Requires: booktabs, tabularx, amsmath, hyperref, biblatex with jtdh_refs.bib

\section*{Appendix: Why Can the LLM Quote Middlemarch but Not Hester?}
\addcontentsline{toc}{section}{Appendix: Analytical Reinforcement}
\label{app:reinforcement}

\noindent
When asked to discuss a novel without source passages, the LLM produces
verifiable quotations at rates ranging from 57\%
(\emph{Middlemarch}) to 0\% (\emph{Hester}).  All 15 novels are
available from Project Gutenberg\footnote{%
  \url{https://www.gutenberg.org}.  Gutenberg texts are a standard
  component of web-scraped training corpora such as C4
  and The Pile.  We assume all 15 novels are present in the training
  data of the model used (Claude, Anthropic), though we cannot verify
  this for any proprietary model.}
and almost certainly appear in the model's training data.  The gradient
must therefore reflect something other than raw textual presence.  This
appendix investigates what that something is, subjects the
investigation to systematic critique, and states what survives.

\subsection*{A.1\quad Hypothesis: Analytical Reinforcement}

\textcite{carlini2023quantifying} showed that LLM memorisation grows
log-linearly with data duplication: sequences that appear more
frequently in training data are reproduced more faithfully.  For
novels, the Gutenberg source text appears once.  But for canonical
novels, specific passages are \emph{duplicated} across thousands of
secondary sources: study guides that excerpt key quotations, academic
papers that discuss them, Wikipedia articles that summarise plots with
embedded excerpts, blog posts, teaching materials, and social media.

A famous passage---the fog opening of \emph{Bleak House}, the pier
glass passage of \emph{Middlemarch}---may appear not only in the
Gutenberg text but across hundreds of such sources.  A passage from
Chapter~12 of \emph{Hester} appears on Gutenberg and virtually
nowhere else.  Our hypothesis: the LLM's per-novel quote fidelity
tracks the volume of this \textbf{analytical reinforcement}, not the
presence of the source text.

\subsection*{A.2\quad Data Sources and Composite Score}

We constructed a composite attention score from three proxies, each
capturing a different facet of a novel's analytical footprint:

\begin{enumerate}[leftmargin=2em]

\item \textbf{Goodreads ratings count.}  Source: Goodreads.com, accessed
March 2026.  This measures present-day readership.
Ratings counts range from 548 (\emph{Hester}) to 261,000
(\emph{David Copperfield}).  Log-normalised to $[0, 100]$ to reduce
the influence of extreme popularity:
$\mathrm{score} = 100 \times (\log r - \log r_{\min}) / (\log r_{\max} - \log r_{\min})$
where $r$ is the ratings count.

\item \textbf{Wikipedia article word count.}  Source: English Wikipedia,
accessed March 2026.  Word count of the novel's Wikipedia article,
measured by selecting all text content.  This measures encyclopaedic
and analytical attention: longer articles contain more plot summary,
critical discussion, and quoted excerpts.  Ranges from $\sim$850
(\emph{Miss Marjoribanks}) to $\sim$19,000 (\emph{David
Copperfield}).  Linearly normalised to $[0, 100]$.

\item \textbf{Study guide presence.}  Binary indicator: does the novel
have a SparkNotes, LitCharts, CliffsNotes, or Shmoop study guide?
Source: manual check of each platform, March 2026.  Study guides are
particularly relevant because they systematically excerpt and re-quote
key passages, producing exactly the kind of passage-level duplication
that \textcite{carlini2023quantifying}'s framework predicts will
strengthen memorisation.  Coded as 33.3 (present) or 0 (absent).
\end{enumerate}

The composite score is the unweighted mean of the three normalised
components, producing a scale of 0--100.  We acknowledge this scoring
is ad hoc: the choice of proxies, normalisation methods, and equal
weighting is not theoretically derived.  Different reasonable choices
would produce different numbers.  We report it as one defensible
operationalisation, not as a validated instrument.

\subsection*{A.3\quad Results}

Table~\ref{tab:attention} presents the full data.

\begin{table}[ht]
\centering
\small
\caption{Critical attention metrics and ungrounded (no-passages) quote
verification rates.  Data sources: Goodreads.com, English Wikipedia,
SparkNotes/LitCharts (all accessed March 2026).  Nop\% is the
percentage of LLM-generated quotes verified against the full enriched
passage corpus using 5-tier matching ($\geq$60\% = verified).}
\label{tab:attention}
\begin{tabularx}{\textwidth}{X|l|r|r|r|c|r|r}
\toprule
\textit{Novel} & \textit{Author} & \textit{Year} &
  \textit{GR ratings} & \textit{Wiki words} & \textit{SG} &
  \textit{Attn.} & \textit{Nop\%} \\
\midrule
Middlemarch        & Eliot    & 1871 & 181,000 & 8,750  & Y &  79.2 & 57 \\
Bleak House        & Dickens  & 1853 &  98,000 & 8,750  & Y &  75.9 & 53 \\
David Copperfield  & Dickens  & 1850 & 261,000 & 19,000 & Y & 100.0 & 44 \\
Hard Times         & Dickens  & 1854 &  75,000 & 6,750  & Y &  70.8 & 39 \\
Passage to India   & Forster  & 1924 &  73,000 & 4,350  & Y &  66.2 & 39 \\
Mill on the Floss  & Eliot    & 1860 &  59,000 & 3,000  & Y &  62.6 & 37 \\
Cranford           & Gaskell  & 1853 &  48,000 & 2,900  & N &  27.9 & 33 \\
Daniel Deronda     & Eliot    & 1876 &  27,000 & 3,500  & Y &  59.3 & 30 \\
Our Mutual Friend  & Dickens  & 1865 &  32,000 & 8,750  & Y &  69.8 & 29 \\
New Grub Street    & Gissing  & 1891 &   6,300 & 3,200  & N &  17.5 & 20 \\
The Odd Women      & Gissing  & 1893 &   5,500 & 1,200  & N &  13.1 &  6 \\
Miss Marjoribanks  & Oliphant & 1866 &   2,600 &   850  & N &   8.4 &  3 \\
North and South    & Gaskell  & 1855 & 182,000 & 8,750  & Y &  79.2 &  2 \\
Hester             & Oliphant & 1883 &     548 & 1,300  & N &   0.8 &  0 \\
No Name            & Collins  & 1862 &   7,900 & 3,200  & N &  18.7 &  0 \\
\bottomrule
\end{tabularx}
\end{table}

\paragraph{Overall correlation.}  Pearson $r = 0.72$, Spearman
$\rho = 0.70$ ($n = 15$).  Excluding \emph{North and South}
(discussed below): Pearson $r = 0.89$, Spearman $\rho = 0.90$
($n = 14$).

\subsection*{A.4\quad The North and South Outlier}

\emph{North and South} has an attention score of 79.2---comparable
to \emph{Middlemarch} and \emph{Bleak House}---but a nop verification
rate of only 2\%.  Its Goodreads count (182,000) is almost certainly
inflated by the 2004 BBC television adaptation, which generated a
large modern fandom that discusses the adaptation's casting, romance,
and divergences from the novel---not the novel's prose passage by
passage.

This suggests that the composite attention score, which treats
Goodreads popularity as equivalent to Wikipedia coverage and study
guide presence, conflates \emph{readership} with \emph{analytical
engagement}.  A novel can be widely read (or widely watched) without
generating the passage-level discussion that reinforces specific
textual fragments in training data.  \emph{Cranford} (also adapted by
the BBC, 2007) shows a milder version of this effect but remains
closer to the trend line.

We report the outlier-excluded correlation ($r = 0.89$) alongside the
full-sample correlation ($r = 0.72$), but we do not rely on the
exclusion.  The more important evidence is the within-author analysis
below.

\subsection*{A.5\quad Within-Author Gradients}

The most damaging objection to the analytical reinforcement hypothesis
is that the correlation might reflect \emph{author effects}: perhaps
Dickens's prose is inherently more memorable (whether by LLMs, by people or by both) than Oliphant's,
regardless of secondary discussion.
We address this by examining
within-author gradients.

\begin{table}[ht]
\centering
\small
\caption{Within-author nop verification gradients.  We assume that novels by the
same author share significant stylistic properties; author effects are mitigated.}
\label{tab:within_author}
\begin{tabularx}{\textwidth}{l|X|r|r}
\toprule
\textit{Author} & \textit{Novel} & \textit{Attn.} & \textit{Nop\%} \\
\midrule
Dickens  & Bleak House        & 75.9 & 53 \\
         & David Copperfield  & 100.0 & 44 \\
         & Hard Times         & 70.8 & 39 \\
         & Our Mutual Friend  & 69.8 & 29 \\
\addlinespace
Eliot    & Middlemarch        & 79.2 & 57 \\
         & Mill on the Floss  & 62.6 & 37 \\
         & Daniel Deronda     & 59.3 & 30 \\
\addlinespace
Gaskell  & Cranford           & 27.9 & 33 \\
         & North and South    & 79.2 &  2 \\
\addlinespace
Gissing  & New Grub Street    & 17.5 & 20 \\
         & The Odd Women      & 13.1 &  6 \\
\addlinespace
Oliphant & Miss Marjoribanks  &  8.4 &  3 \\
         & Hester             &  0.8 &  0 \\
\bottomrule
\end{tabularx}
\end{table}

\textbf{Dickens} (4 novels): Bleak House (53\%) $>$ David Copperfield
(44\%) $>$ Hard Times (39\%) $>$ Our Mutual Friend (29\%).  This
within-Dickens ranking tracks canonical status: Bleak House is the
most taught and discussed; Our Mutual Friend, though artistically
admired, is the least widely assigned.  Same author, plausibly similar prose style,
different critical prominence, different quote fidelity.

\textbf{Eliot} (3 novels): Middlemarch (57\%) $>$ Mill on the Floss
(37\%) $>$ Daniel Deronda (30\%). It is plausible that Middlemarch is far more widely
taught and discussed than Daniel Deronda.

\textbf{Gaskell} (2 novels): Cranford (33\%) $\gg$ North and South
(2\%).  This is the most striking within-author contrast and the one
we understand least well.  North and South has more study guide
coverage (GradeSaver, LitCharts, Course Hero, SuperSummary) and a
higher attention score (79.2 vs.\ 27.9) than Cranford, yet Cranford
dramatically outperforms.  North and South's modern Goodreads
popularity is inflated by the 2004 BBC adaptation, which may generate
discussion of the \emph{adaptation} without reinforcing the
\emph{novel's} specific passages.  But we cannot rule out other
explanations: Cranford's shorter length may make it more quotable per
passage, or its comic tone may produce more distinctive (and therefore
more memorisable) phrasing.  This pair is genuinely anomalous and
resists the simple analytical-reinforcement explanation.

\textbf{Gissing} (2 novels): New Grub Street (20\%) $>$ The Odd Women
(6\%).  New Grub Street is the better-known novel.

\textbf{Oliphant} (2 novels): Miss Marjoribanks (3\%) $\approx$
Hester (0\%).  Both near zero, consistent with Oliphant's general
absence from modern critical discussion.

\medskip\noindent
Within-author gradients exist for all five authors with multiple novels
and, with the exception of the Gaskell pair (discussed in A.4), they
track critical-cultural prominence.  \textbf{This is the strongest
evidence for the analytical reinforcement hypothesis}, because it
controls for the most plausible confound: author prose style.

\subsection*{A.6\quad Objections and Qualifications}

We subjected the thesis to systematic critique.  The principal
objections, and our responses:

\begin{description}[leftmargin=1.5em, labelindent=0em]

\item[Training data opacity.]  We do not know the exact composition of
the model's training data.  The attention score is a proxy for likely
duplication in web-scraped corpora, not a measurement of actual
duplication.  The hypothesis is consistent with
\textcite{carlini2023quantifying}'s findings but cannot be verified
for proprietary models.

\item[Ad hoc scoring.]  The composite score is one defensible
operationalisation among several.  We tested robustness informally:
dropping any single proxy changes $r$ by less than 0.1.  A more
principled instrument would require validated cultural-prominence
metrics that do not currently exist.

\item[Small sample.]  $n = 15$ supports trend identification but not
confident statistical inference.  The correlation ($r = 0.72$) has a
wide confidence interval.  We rely on the within-author evidence
(Table~\ref{tab:within_author}) as much as on the aggregate
correlation.

\item[Readership vs.\ analytical engagement.]  The North and South
outlier demonstrates that popular readership does not entail the kind
of passage-level analytical engagement that would reinforce specific
fragments in training data.  A more refined metric would distinguish
passage-quoting activity (study guides, academic close readings) from
generic discussion (reviews, adaptation commentary).  We have not
measured this distinction directly.

\item[The Bayard parallel is a metaphor.]
\textcite{bayard2007comment}'s argument concerns human cultural
navigation; our findings concern statistical weight distributions in a
language model.  The parallel---both humans and LLMs navigate literary
culture through graduated approximation---is evocative but the
mechanisms differ entirely.  We flag it as an analogy, not an identity.
As \textcite{shanahan2024talking} argues, care is needed in mapping
LLM behaviour onto human cognitive categories.

\end{description}

\subsection*{A.7\quad Conclusions}

\textbf{Stated conservatively:}

All 15 novels are available from Project Gutenberg and almost certainly
appear in LLM training data.  Yet the LLM's ability to quote them from
memory varies from 57\% to 0\%.  This gradient correlates with
critical-cultural prominence ($r = 0.72$, $n = 15$) and persists
within authors: among four Dickens novels, three Eliot novels, and two
Gissing novels, the within-author ranking tracks the novel's position
in modern critical culture.

We interpret this as a literary-critical manifestation of differential
LLM memorisation: canonical novels, whose passages are duplicated
across study guides, academic papers, and Wikipedia, are more
faithfully retained than obscure ones.  This interpretation is
consistent with \textcite{carlini2023quantifying}'s framework but
unverifiable without access to training data.

The parallel with \textcite{bayard2007comment}'s argument about human
non-reading is suggestive: both humans and LLMs navigate literary
culture through graduated approximation.  The mechanisms are entirely
different, but the observable gradient---from faithful quotation
of canonical works to creative pastiche of obscure ones---is
strikingly similar.
